% Generated draft: requires mathematical and source review.
\documentclass[11pt]{article}
\usepackage[margin=25mm]{geometry}
\usepackage{fontspec}
\setmainfont{Latin Modern Roman}
\usepackage{amsmath,amssymb,mathtools,unicode-math}
\setmathfont{Latin Modern Math}
\usepackage{enumitem,needspace,xurl,hyperref}
\setlength{\parindent}{0pt}
\setlength{\parskip}{5pt}
\setlength{\emergencystretch}{3em}
\newcommand{\N}{\mathbb{N}}
\newcommand{\Z}{\mathbb{Z}}
\newcommand{\Q}{\mathbb{Q}}
\newcommand{\R}{\mathbb{R}}
\newcommand{\C}{\mathbb{C}}
\newcommand{\F}{\mathbb{F}}
\newcommand{\A}{\mathbb{A}}
\newcommand{\PP}{\mathbb{P}}
\newcommand{\E}{\mathbb{E}}
\newcommand{\eps}{\varepsilon}
\newcommand{\dd}{\,\mathrm{d}}
\begin{document}
\section*{UnsolvedMath ID: 30005861}
\textbf{Source problem code:} catalogue entry
\subsection{Definitions and mathematical statement}
The supplied record combines Open Questions 5 and 6 from the cited workshop report. They are requests for geometric hypotheses and a convergence framework; the supplied excerpt does not specify enough data to give a single fully quantified conjecture.

For a smooth surface \(\Sigma\) in a setting admitting a chosen unit normal, the constant-mean-curvature condition is
\[
H_\Sigma\equiv c
\]
for some constant \(c\in\R\). Given a Riemannian manifold \((M,g)\), an admissible class of regions, and a perimeter convention \(P_g\), an isoperimetric region of prescribed volume \(v\) is an admissible region \(\Omega\) satisfying
\[
\operatorname{Vol}_g(\Omega)=v,
\qquad
P_g(\Omega)=\inf\{P_g(E): E\text{ admissible},\ \operatorname{Vol}_g(E)=v\}.
\]
The admissible volumes, boundary convention, and existence or regularity assumptions are not specified in the record.

\textbf{Open Question 5: geometric noncollapse.} Formulate noncollapsing hypotheses for sequences of Riemannian manifolds using constant-mean-curvature surfaces or isoperimetric regions. The source does not specify a quantitative meaning of noncollapse, the geometric quantities to control, or the conclusion those hypotheses should imply. Consequently, no particular volume lower bound or other noncollapse criterion is asserted here.

\textbf{Open Question 6: convergence retaining Sobolev information.} For a sequence of varying Riemannian manifolds \((M_j,g_j)\), construct a notion of convergence that implies volume-preserving intrinsic-flat convergence and also captures information encoded in a \(W^{1,p}\) limit. The supplied reconstructed statement describes retaining this information on ``good regions.'' Determine additionally whether a distance between Riemannian manifolds can capture the proposed convergence.

To make the intrinsic-flat requirement explicit, suppose the manifolds and a candidate limit \((M_\infty,g_\infty)\) determine integral current spaces \(X_j\) and \(X_\infty\), respectively, in a setting where intrinsic-flat distance \(d_{\mathcal F}\) and finite total Riemannian volume are defined. In this setting, the intended volume-preserving requirement is
\[
d_{\mathcal F}(X_j,X_\infty)\longrightarrow 0,
\qquad
\operatorname{Vol}_{g_j}(M_j)\longrightarrow
\operatorname{Vol}_{g_\infty}(M_\infty).
\]
Here intrinsic-flat distance is defined by comparing the pushed-forward integral currents in common complete metric spaces using flat distance. This notation describes the requirement when the limit is represented by a Riemannian manifold; the source does not specify whether more general limit spaces are intended.

The notation \(W^{1,p}\) refers to Sobolev regularity: the relevant functions or tensor coefficients and their first weak derivatives have the prescribed \(L^p\) integrability in chosen coordinates or identifications. On varying manifolds, a precise convergence assertion additionally requires identification maps, specified regions, specified objects to compare, an exponent \(p\), and a choice of strong or weak convergence. None of these is supplied. In particular, ``good regions'' has no supplied definition. Likewise, the class of manifolds and the equivalence relation on which the requested distance should be a metric remain unspecified.

The report title supplies a positive-scalar-curvature context, but the excerpt gives no precise scalar-curvature hypothesis, dimension, compactness assumption, boundary condition, or other uniform geometric bounds. Such hypotheses must be recovered from the surrounding source before either question can be made fully precise.
\subsection{Short English statement}
The record asks for noncollapse conditions expressed through constant-mean-curvature surfaces or isoperimetric regions, and for convergence of varying Riemannian manifolds that preserves volume under intrinsic-flat convergence while retaining \(W^{1,p}\) information. It also asks whether a distance can describe that convergence. The geometric setting, Sobolev comparisons, and meaning of ``good regions'' remain unspecified, so this is a qualified research-program statement requiring human mathematical review.
\subsection{Sources}
\begin{itemize}
\item Local UnsolvedMath registry: permanent ID \(30005861\), title \emph{Geometric Noncollapse and Hybrid Convergence}, source problem code \texttt{catalogue entry}, in the Oberwolfach Reports Open Problems collection. The supplied record combines Open Questions 5 and 6 and lists the proposed year as 2024; no individual proposer is identified.
\item \emph{Analysis, Geometry and Topology of Positive Scalar Curvature Metrics} (2024), Oberwolfach Reports, pp.~483--566, DOI 10.4171/owr/2024/9. The supplied excerpt identifies p.~549. Original report: \url{https://doi.org/10.4171/owr/2024/9}.
\item Supplied literature source, labeled ``Sormani--Tian--Wang scalar-curvature limit report'': \url{https://arxiv.org/abs/2404.17121}.
\item Supplied literature source, labeled ``Allen scalar-curvature stability report'': \url{https://arxiv.org/abs/2404.02662}.
\item Literature assessment checked at \texttt{2026-08-22}: the supplied dated curation labels the problem open and reports no established general construction or equivalence theorem simultaneously encoding volume-preserving SWIF convergence and \(W^{1,p}\) information in the literature checked. This is supplied dated curation, not an independently verified assessment; the record requests independent verification of time-sensitive status and cited claims.
\item The registry labels its statement \texttt{corrected\_verified} and reports that it was reconstructed and checked against the official Oberwolfach report. These are supplied provenance assertions; no source visit or independent mathematical verification is claimed here.
\end{itemize}
\end{document}
